\documentclass[11pt,letterpaper]{article}
% Standalone combined source: all content and formatting are embedded.
% Compact questions/solutions layout matching Versions A--C.
\usepackage[margin=1in]{geometry}
\usepackage[T1]{fontenc}
\usepackage{amsmath,amssymb}
\usepackage{enumitem}
\usepackage{needspace}
\usepackage{xurl}
\usepackage[hidelinks]{hyperref}
\hypersetup{hypertexnames=false}

\setlength{\parindent}{0pt}
\setlength{\parskip}{4pt}
\setlength{\emergencystretch}{1.5em}
\setlength{\abovedisplayskip}{5pt}
\setlength{\belowdisplayskip}{5pt}
\setlength{\abovedisplayshortskip}{3pt}
\setlength{\belowdisplayshortskip}{3pt}
\setlist[itemize]{leftmargin=1.4em,itemsep=2pt,topsep=3pt}
\setlist[enumerate]{leftmargin=1.7em,itemsep=4pt,topsep=4pt}
\newcommand{\curl}{\nabla\times}
\newcommand{\dv}{\nabla\cdot}
\renewcommand{\vec}[1]{\mathbf{#1}}
\newcommand{\dd}{\mathop{}\!\mathrm{d}}
\newcommand{\grad}{\nabla}
\allowdisplaybreaks[2]
\pagestyle{plain}

\newcommand{\examproblem}[3]{%
  \par\medskip\Needspace{5\baselineskip}
  \noindent\textbf{#1. #2}\par\nobreak
}
\newcommand{\solutionproblem}[3]{%
  \par\bigskip\Needspace{7\baselineskip}
  \noindent\textbf{#1. #2}\par\nobreak
}
\newcommand{\insight}[1]{\par\textbf{Key observation.} #1\par}
\newcommand{\verification}[1]{\par\textbf{Check.} #1\par}
\newcommand{\alternative}[1]{\par\textbf{Alternative.} #1\par}

\begin{document}
{\large\bfseries Mock Exam 3 --- Version A --- Solutions}
\smallskip
% Worked solutions; provenance notes are retained in comments below.
\solutionproblem{1}{Curl, divergence, and conservativity}{10}
\insight{The $\sin z$ and $\cos z$ terms come from the gradient of
$xy\sin z$; the remaining planar rotation supplies a nonzero curl.}
Writing $\vec F=(P,Q,R)$ gives
\[
\begin{aligned}
R_y-Q_z&=x\cos z-x\cos z=0,\\
P_z-R_x&=y\cos z-y\cos z=0,\\
Q_x-P_y&=(\sin z+1)-(\sin z-1)=2.
\end{aligned}
\]
Therefore
\[
\boxed{\curl\vec F=(0,0,2)},\qquad
\boxed{\dv\vec F=-xy\sin z}.
\]
The divergence follows from $P_x=Q_y=0$ and $R_z=-xy\sin z$.
The field is \textbf{not conservative}: a gradient field with continuous
second partial derivatives has zero curl, whereas this curl is nonzero
everywhere.
\verification{The decomposition
$\vec F=\nabla(xy\sin z)+(-y,x,0)$ independently gives the curl
$(0,0,2)$ and divergence $-xy\sin z$.}
\alternative{Computing all six mixed-component derivatives directly is
fully sufficient; recognizing the decomposition is optional.}

\solutionproblem{2}{Green's theorem}{12}
\insight{The exponential terms have equal mixed derivatives. After
they cancel, the integral depends only on areas and a first moment of
an off-center hole. The outer disk's odd moment is zero, but the hole's
moment is not.}
Let $P=ye^{xy}-y$ and $Q=xe^{xy}+3x^2+x$. Then
\[
Q_x=e^{xy}+xye^{xy}+6x+1,\qquad
P_y=e^{xy}+xye^{xy}-1,
\]
so $Q_x-P_y=6x+2$. The given orientation is positive: the outer circle
is counterclockwise and the inner circle clockwise. Green's theorem gives
\[
\oint_C P\,\dd x+Q\,\dd y=\iint_D(6x+2)\,\dd A.
\]
The smaller disk lies strictly inside the larger one because
$1+1/2<2$. An explicit setup is
\[
\begin{aligned}
\iint_D(6x+2)\,\dd A
={}&\int_0^{2\pi}\int_0^2(6r\cos\theta+2)r\,\dd r\,\dd\theta\\
&-\int_0^{2\pi}\int_0^{1/2}
\bigl(6(1+s\cos\theta)+2\bigr)s\,\dd s\,\dd\theta.
\end{aligned}
\]
The second integral uses coordinates centered at $(1,0)$, not at the
origin. Terms containing $\cos\theta$ integrate to zero. Thus
\[
\oint_C P\,\dd x+Q\,\dd y
=2(4\pi)-8\left(\frac\pi4\right)
=\boxed{6\pi}.
\]
\verification{The area of $D$ is $15\pi/4$ and
$\iint_D x\,\dd A=0-(1)(\pi/4)=-\pi/4$; hence
$6(-\pi/4)+2(15\pi/4)=6\pi$.}
\alternative{Cartesian bounds, splitting the vertical slices through
the hole, are valid. Translating the small disk and using its symmetry
avoids that split. A proof of the first-moment calculation by integrating
in the displayed translated coordinates is already given above.}

\solutionproblem{3}{Three surface parametrizations}{12}
\insight{Use coordinates suited to the geometry, then translate every
extra condition into a parameter bound. No surface-area calculation is
needed.}
\begin{enumerate}[label=(\alph*)]
\item In polar coordinates the cylinder condition becomes
\[
(r\cos\theta-1)^2+r^2\sin^2\theta\le1
\quad\Longleftrightarrow\quad r(r-2\cos\theta)\le0.
\]
Thus one parametrization is
\[
\boxed{\vec r(r,\theta)=(r\cos\theta,r\sin\theta,r)},\qquad
\left\{\begin{aligned}
&-\pi/2\le\theta\le\pi/2,\\
&0\le r\le2\cos\theta.
\end{aligned}\right.
\]
The cone's tip is a boundary point; degeneracy of polar coordinates
there does not affect the image.

\item Scale spherical coordinates to the ellipsoid:
\[
\boxed{\vec r(\phi,\theta)
=(2\sin\phi\cos\theta,3\sin\phi\sin\theta,\cos\phi)},\qquad
\left\{\begin{aligned}
&0\le\phi\le\pi/4,\\
&0\le\theta\le\pi.
\end{aligned}\right.
\]
The cone-like inequality is $\cos\phi\ge\sin\phi$, so it forces
$0\le\phi\le\pi/4$; $y\ge0$ then gives the stated azimuthal range.

\item A parametrization using the taught $y$-axis revolution convention is
\[
\boxed{\vec r(s,\theta)
=\bigl((2+\cos s)\cos\theta,\sin s,-(2+\cos s)\sin\theta\bigr)},
\]
\[
\left\{\begin{aligned}
&-\pi/2\le s\le\pi/2,\\
&0\le\theta\le2\pi.
\end{aligned}\right.
\]
Here the distance to the $y$-axis is $2+\cos s\ge2$, so this is an
ordinary surface of revolution without a crossing of the axis.
\end{enumerate}
\verification{In (a), $z=r=\sqrt{x^2+y^2}$. In (b), substitution gives
$x^2/4+y^2/9+z^2=1$ exactly. In (c), setting $\theta=0$ recovers the
given generating arc, and varying $\theta$ preserves its height and
distance from the axis.}
\alternative{Any two-parameter descriptions with the same images and
correct complete bounds are valid. Reversing angular direction, shifting
angles by $2\pi$, using a different pair of parameter names, and different
once-around representations are equivalent. No orientation is requested.}

\solutionproblem{4}{A scalar surface integral: setup only}{12}
\insight{On the supplied surface, $x^2-y^2=4z$, so the exponential
collapses to 1 before the scalar integrand is multiplied by the area
factor.}
The given parametrization has
\[
\vec r_u=(1,1,v),\qquad \vec r_v=(1,-1,u),
\]
\[
\vec r_u\times\vec r_v=(u+v,v-u,-2),\qquad
\lVert\vec r_u\times\vec r_v\rVert=\sqrt{4+2u^2+2v^2}.
\]
The scalar substitution is
\[
f(\vec r(u,v))
=e^{(u+v)^2-(u-v)^2-4uv}+uv
=1+uv.
\]
Consequently the requested setup is
\[
\boxed{\iint_S f\,\dd\sigma
=\int_0^1\int_0^{1-u}(1+uv)\sqrt{4+2u^2+2v^2}\,\dd v\,\dd u}.
\]
\verification{The squared cross-product norm is
$(u+v)^2+(v-u)^2+4=4+2u^2+2v^2>0$. The first two surface coordinates
uniquely recover $u=(x+y)/2$ and $v=(x-y)/2$, so the patch is not
multiply covered. The density $1+uv$ is nonnegative on the parameter
triangle.}
\alternative{Reversing the integration order gives bounds
$0\le v\le1$, $0\le u\le1-v$ with the same scalar integrand. The
opposite cross-product order is valid because its magnitude is unchanged.}

\solutionproblem{5}{Direct flux through a curved surface}{20}
\insight{The terms involving $e^{z^2}$ describe motion around the axis,
so their normal contribution cancels. A second cancellation uses
$z=x^2+y^2$; the surviving integral has the derivative needed for
$e^{r^2}$.}
Choose
\[
\vec r(r,\theta)=(r\cos\theta,r\sin\theta,r^2),\qquad
\left\{\begin{aligned}
&1\le r\le2,\\
&0\le\theta\le2\pi.
\end{aligned}\right.
\]
The tangent vectors and outward cross product are
\[
\vec r_r=(\cos\theta,\sin\theta,2r),\qquad
\vec r_\theta=(-r\sin\theta,r\cos\theta,0),
\]
\[
\vec n=\vec r_\theta\times\vec r_r
=(2r^2\cos\theta,2r^2\sin\theta,-r).
\]
At $(r,\theta)=(1,0)$ this gives $(2,0,-1)$, whose radial component
points away from the $z$-axis. More generally,
$\vec n\cdot(r\cos\theta,r\sin\theta,0)=2r^3>0$, verifying the
required orientation.

Set $c=\cos\theta$, $s=\sin\theta$. Substitution gives
\[
\vec F(\vec r)=
\bigl(rc-rs e^{r^4},\ rs+rc e^{r^4},\ 2r^2-e^{r^2}\bigr).
\]
The dot product simplifies before any integration:
\[
\begin{aligned}
\vec F(\vec r)\cdot\vec n
&=2r^3c^2-2r^3sc e^{r^4}
 +2r^3s^2+2r^3sc e^{r^4}
 -2r^3+r e^{r^2}\\
&=r e^{r^2}.
\end{aligned}
\]
Thus
\[
\boxed{\iint_S\vec F\cdot\vec N\,\dd\sigma
=\int_0^{2\pi}\int_1^2r e^{r^2}\,\dd r\,\dd\theta
=\pi(e^4-e)}.
\]
\verification{The vector $(x,y,2z)$ is tangent to $z=x^2+y^2$ because
its dot product with $(2x,2y,-1)$ is zero on that surface.
The rotation vector $(-y,x,0)$ is also tangent. Only $(0,0,-e^z)$
contributes to the outward flux, and its downward component agrees with
the downward component of the outward normal, yielding positive flux.}
\alternative{A direct graph parametrization $(x,y,x^2+y^2)$ with
outward vector $(2x,2y,-1)$ produces integrand $e^{x^2+y^2}$ over
$1\le x^2+y^2\le4$. Converting that double integral to polar gives
the same result. Both methods directly parametrize the stated surface.}

\solutionproblem{6}{Stokes' theorem and the whole boundary}{20}
\insight{The lateral cone has two boundary circles, at different heights.
Their induced directions are opposite, and the coefficient $z+1$ takes
a different value on each. The third component of the field contributes
nothing because both circles have $\dd z=0$.}
An outward normal on the upper cone is proportional to
$(\cos\theta,\sin\theta,-1)$, so it points downward. Consequently the
outer circle at $z=3$ is clockwise from above, and the inner circle at
$z=1$ is counterclockwise.

One-parameter parametrizations with the induced orientations are
\[
\begin{aligned}
\vec c_o(t)&=(3\cos t,-3\sin t,3),\\
\vec c_i(t)&=(\cos t,\sin t,1),
\end{aligned}
\qquad0\le t\le2\pi.
\]
Their derivatives and substituted fields are
\[
\begin{aligned}
\vec c_o'(t)&=(-3\sin t,-3\cos t,0),&
\vec F(\vec c_o(t))&=(12\sin t,12\cos t,e^{27}),\\
\vec c_i'(t)&=(-\sin t,\cos t,0),&
\vec F(\vec c_i(t))&=(-2\sin t,2\cos t,e).
\end{aligned}
\]
Thus
\[
\vec F(\vec c_o(t))\cdot\vec c_o'(t)=-36,\qquad
\vec F(\vec c_i(t))\cdot\vec c_i'(t)=2.
\]
Stokes' theorem, with the whole oriented boundary, gives
\[
\begin{aligned}
\iint_S(\curl\vec F)\cdot\vec N\,\dd\sigma
&=\oint_{C_o}\vec F\cdot\dd\vec r
 +\oint_{C_i}\vec F\cdot\dd\vec r\\
&=\int_0^{2\pi}(-36)\,\dd t
 +\int_0^{2\pi}2\,\dd t
=\boxed{-68\pi}.
\end{aligned}
\]
\verification{On a counterclockwise circle of radius $R$ at height $h$,
the first two field components give circulation
$2\pi(h+1)R^2$. Reversing the outer circle therefore gives
$-2\pi(4)(9)+2\pi(2)(1)=-68\pi$, confirming both signs and radii.}
\alternative{Using counterclockwise parametrizations for both circles
and subtracting the outer-circle integral is equivalent. A direct
surface calculation also verifies the number independently.}

\solutionproblem{7}{Conceptual decisions}{14}
\begin{enumerate}[label=(\alph*)]
\item \[
\begin{array}{c|c|c}
\text{Expression}&\text{Type}&\text{Necessarily zero?}\\\hline
\curl(f\vec G)&\text{vector field}&\text{No}\\
\dv(\curl\vec G)&\text{scalar field}&\text{Yes}\\
\curl(\dv\vec G)&\text{undefined}&\text{not applicable}
\end{array}
\]
For the first row, $f=1$ and $\vec G=(-y,x,0)$ give curl $(0,0,2)$,
so a universal zero claim is false. The second row is the identity
$\dv(\curl\vec G)=0$. In the third row, divergence produces a scalar,
but curl takes a vector field as its input.

\item The only directly applicable listed method is
\textbf{direct parametrization}. The curve is open, so Green's and
Stokes' theorems do not directly apply. Although
\[
\curl(0,xy,0)=(0,0,y)
\]
vanishes along $C$, it is nonzero at points with $y\ne0$, including
points arbitrarily close to the curve. Hence the stated field is not
conservative on an open neighborhood of $C$, and FTLI cannot be invoked
for it. The fact that its particular line integral happens to equal zero
does not supply the missing FTLI hypothesis.

\item The negative divergence indicates a local
\textbf{sink}, or local net inflow. The directed rotation axis is
\textbf{positive $z$}, and the right-hand rule gives
\textbf{counterclockwise} rotation viewed from the positive $z$-axis.
These are local conclusions at $p$, not global statements about the flow.
\end{enumerate}
\verification{For (b), direct substitution gives $\vec F(t,0,0)=\vec0$,
so direct evaluation is possible even though FTLI is not applicable to
the given nonconservative field. For (c), the signs and axis are read
from the supplied values, without needing a field formula.}
\alternative{Any equivalent physical wording for net local inflow and
the positive-$z$ right-hand-rule sense is acceptable. No global behavior
or trajectory shape is determined by those two local measurements.}

\clearpage
\setcounter{page}{1}
{\large\bfseries Mock Exam 3 --- Version B --- Solutions}
\smallskip
% Worked solutions; provenance notes are retained in comments below.
\solutionproblem{1}{Curl, divergence, and conservativity}{10}
\insight{The exponential and trigonometric pieces have matching mixed
partials. Their cancellations must be performed before deciding whether
the remaining curl vanishes everywhere.}
Writing $\vec F=(P,Q,R)$ gives
\[
\begin{aligned}
R_y-Q_z&=2z-(2z+x)=-x,\\
P_z-R_x&=(\cos x-y)-\cos x=-y,\\
Q_x-P_y&=(e^{xy}+xye^{xy}+z)
 -(e^{xy}+xye^{xy}-z)=2z.
\end{aligned}
\]
Consequently
\[
\boxed{\curl\vec F=(-x,-y,2z)},\qquad
\boxed{\dv\vec F=(x^2+y^2)e^{xy}-z\sin x+2y}.
\]
The field is \textbf{not conservative} on $\mathbb R^3$, since its curl
is not identically zero; for example, it equals $(-1,0,0)$ at $(1,0,0)$.
Its zero curl at the origin alone cannot establish conservativity.
\verification{The first three terms form
$\grad(e^{xy}+z\sin x+yz^2)$. The remaining field
$(-yz,xz,0)$ has curl $(-x,-y,2z)$ and divergence zero, independently
confirming both the cancellation structure and the final formulas.}
\alternative{Differentiate the three components directly, as above. A
potential need not be found because the nonzero curl already rules one
out.}

\solutionproblem{2}{Green's theorem on an unfinished curve}{12}
\insight{Close the missing diagonal, use the clockwise orientation of the
resulting curve, and subtract the nonzero integral on that diagonal.}
Let $L$ be the segment from $(2,1)$ back to $(1,0)$. The closed curve
$C+L$ traverses its region clockwise. The enclosed region is
\[
D:\qquad
\left\{
\begin{aligned}
&0\le y\le1,\\
&y^2\le x\le1+y.
\end{aligned}
\right.
\]
For $P=ye^{xy}-y^2/2$ and $Q=xe^{xy}+2x$,
\[
Q_x-P_y=(e^{xy}+xye^{xy}+2)
 -(e^{xy}+xye^{xy}-y)=2+y.
\]
Green's theorem, with the clockwise sign, gives
\[
\begin{aligned}
\int_C\vec F\cdot\dd\vec r+\int_L\vec F\cdot\dd\vec r
&=-\int_0^1\int_{y^2}^{1+y}(2+y)\,\dd x\,\dd y\\
&=-\int_0^1(2+3y-y^2-y^3)\,\dd y=-\frac{35}{12}.
\end{aligned}
\]
Parametrize the closing segment by
$\vec\ell(t)=(1+t,t)$, with $t$ decreasing from $1$ to $0$. Then
\[
\begin{aligned}
\int_L\vec F\cdot\dd\vec r
&=\int_1^0\left((1+2t)e^{t+t^2}-\frac{t^2}{2}+2+2t\right)\dd t\\
&=\left[e^{t+t^2}-\frac{t^3}{6}+2t+t^2\right]_1^0
=-e^2-\frac{11}{6}.
\end{aligned}
\]
Therefore
\[
\boxed{\displaystyle\int_C\vec F\cdot\dd\vec r
=-\frac{35}{12}-\left(-e^2-\frac{11}{6}\right)
=e^2-\frac{13}{12}}.
\]
\verification{For an independent check, the original three pieces have
integrals $0$, $e-7/12$, and $e^2-e-1/2$. Their sum agrees with the
answer. This directly confirms the result.}
\alternative{One may close by an easier two-segment polygonal route,
provided every auxiliary integral and the resulting oriented region are
accounted for. The diagonal gives a single closing integral.}

\solutionproblem{3}{Surface parametrizations and intersection bounds}{12}
\insight{The equations determine parameter bounds by intersections, not
just by identifying a familiar surface.}
\textbf{(a)} One parametrization is
\[
\boxed{\vec r(\theta,\phi)
=(2\sin\phi\cos\theta,2\sin\phi\sin\theta,2\cos\phi)},
\qquad
\boxed{0\le\theta\le2\pi,\quad\pi/6\le\phi\le\pi/2}.
\]
The outside-cylinder condition is $4\sin^2\phi\ge1$. Since $z\ge0$,
$0\le\phi\le\pi/2$, so this becomes $\phi\ge\pi/6$.

\textbf{(b)} A convenient parametrization and bounds are
\[
\boxed{\vec r(\theta,v)=(\cos\theta,\sin\theta,v)},\qquad
\boxed{0\le\theta\le2\pi,\quad\sin\theta\le v\le2+\cos\theta}.
\]
The upper plane is always above the lower one because
$2+\cos\theta-\sin\theta\ge2-\sqrt2>0$.

\textbf{(c)} Using the radius as a parameter,
\[
\boxed{\vec r(u,\theta)=(u\cos\theta,u\sin\theta,u)},\qquad
\boxed{1\le u\le\sqrt2,\quad0\le\theta\le2\pi}.
\]
On the cone, the sphere condition becomes $2u^2\le4$, while $z\ge1$
becomes $u\ge1$.
\verification{Substitution gives the sphere, cylinder, and cone equations
respectively. Their bounds include each required limiting intersection
and add no caps.}
\alternative{In (b), one can instead use
\[
\vec r(\theta,t)=(\cos\theta,\sin\theta,
\sin\theta+t(2+\cos\theta-\sin\theta)),
\]
with
$0\le\theta\le2\pi$, $0\le t\le1$. Other correct two-parameter maps
with equivalent domains are valid.}

\solutionproblem{4}{Set up a scalar surface integral}{12}
\insight{The given coordinates make $y-x^2=v$ and $z/x=v$. The latter
identity is legitimate because $x=u\ge1$ everywhere on the patch.}
The tangent vectors and cross product are
\[
\begin{aligned}
\vec r_u&=(1,2u,v),\qquad \vec r_v=(0,1,u),\\
\vec r_u\times\vec r_v&=(2u^2-v,-u,1),\\
\left|\vec r_u\times\vec r_v\right|
&=\sqrt{(2u^2-v)^2+u^2+1}.
\end{aligned}
\]
Substitution gives
\[
f(\vec r(u,v))=e^{(u^2+v)-u^2}
+\left(\frac{uv}{u}\right)^2=e^v+v^2.
\]
Thus the requested setup is
\[
\boxed{\displaystyle\int_1^2\int_0^{2-u}
(e^v+v^2)\sqrt{(2u^2-v)^2+u^2+1}\,\dd v\,\dd u}.
\]
It should not be evaluated.
\verification{The third component of the cross product is $1$, so the
patch is regular throughout its domain. Also $u=x$ and $v=y-x^2$
recover the parameters uniquely; the map does not cover the surface
multiple times.}
\alternative{Reversing the cross-product order gives the same magnitude.
The equivalent bounds $0\le v\le1$, $1\le u\le2-v$ are also valid.}

\solutionproblem{5}{Flux through a curved sheet}{20}
\insight{The first-quadrant lens is unchanged by exchanging $x$ and $y$.
After forming the correctly oriented flux integrand, its antisymmetric
part integrates to zero.}
The two boundary curves meet at $(0,0)$ and $(1,1)$. The region has bounds
\[
D:\qquad
\left\{
\begin{aligned}
&0\le x\le1,\\
&x^2\le y\le\sqrt x.
\end{aligned}
\right.
\]
Use $\vec r(x,y)=(x,y,xy)$. Then
\[
\vec r_x=(1,0,y),\qquad \vec r_y=(0,1,x),\qquad
\vec r_x\times\vec r_y=(-y,-x,1).
\]
This cross product points upward. The required downward oriented area
vector is therefore $\vec r_y\times\vec r_x=(y,x,-1)$; its third
component is negative at every point.

The field and dot product become
\[
\begin{aligned}
\vec F(\vec r)&=(x+1,y-1,3xy+x-y),\\
\vec F(\vec r)\cdot(y,x,-1)
&=(x+1)y+(y-1)x-(3xy+x-y)\\
&=-xy-2x+2y.
\end{aligned}
\]
The complete flux setup is
\[
\iint_S\vec F\cdot\vec N\,\dd\sigma
=\int_0^1\int_{x^2}^{\sqrt x}(-xy-2x+2y)\,\dd y\,\dd x.
\]
Because $D$ is symmetric under $x\leftrightarrow y$,
$\iint_D(-2x+2y)\,\dd A=0$. Hence
\[
\begin{aligned}
\iint_S\vec F\cdot\vec N\,\dd\sigma
&=-\int_0^1\int_{x^2}^{\sqrt x}xy\,\dd y\,\dd x\\
&=-\frac12\int_0^1(x^2-x^5)\,\dd x
=-\frac12\left(\frac13-\frac16\right)
=\boxed{-\frac1{12}}.
\end{aligned}
\]
\verification{Direct integration gives
$\iint_D x\,\dd A=\iint_D y\,\dd A=3/20$ and
$\iint_D xy\,\dd A=1/12$, confirming the cancellation and sign.}
\alternative{Evaluate the displayed double integral without symmetry, or
reverse the order using $0\le y\le1$, $y^2\le x\le\sqrt y$. Both
retain the direct surface-parametrization method.}

\solutionproblem{6}{Stokes' theorem with a nonplanar boundary}{20}
\insight{On the boundary $x^2+y^2=1$, two complicated field terms vanish.
The remaining third component becomes a function of $z$ alone, so its
line-integral contribution is a total derivative around a closed curve.}
The boundary is
\[
C:\qquad x^2+y^2=1,\qquad z=x^2-y^2.
\]
The upward orientation induces counterclockwise traversal of its
$xy$-projection. Thus a matching one-parameter parametrization is
\[
\vec c(t)=(\cos t,\sin t,\cos2t),\qquad0\le t\le2\pi,
\]
with
\[
\vec c\,'(t)=(-\sin t,\cos t,-2\sin2t).
\]
Stokes' theorem gives
\[
\iint_S(\curl\vec F)\cdot\vec N\,\dd\sigma
=\int_0^{2\pi}\vec F(\vec c(t))\cdot\vec c\,'(t)\,\dd t.
\]
Put $z(t)=\cos2t$. Since $1-x^2-y^2=0$ on $C$,
\[
\vec F(\vec c(t))
=\bigl(-\sin t\,z(t)^2,\ \cos t\,z(t)^2,\ e+z(t)e^{z(t)}\bigr).
\]
Consequently the dot product is
\[
z(t)^2+(e+z(t)e^{z(t)})z'(t)
=\cos^2 2t-2\sin2t\left(e+\cos2t\,e^{\cos2t}\right).
\]
The primitive
\[
\Phi(z)=ez+(z-1)e^z
\quad\text{satisfies}\quad
\Phi'(z)=e+ze^z.
\]
Since $z(0)=z(2\pi)=1$, the second term integrates to zero. Therefore
\[
\boxed{\displaystyle
\iint_S(\curl\vec F)\cdot\vec N\,\dd\sigma
=\int_0^{2\pi}\cos^2 2t\,\dd t
 +[\Phi(z(t))]_0^{2\pi}=\pi}.
\]
\verification{The terms multiplied by $1-x^2-y^2$ are exactly zero on
every boundary point. For the remaining field, a separate surface check
has integrand $6(x^2-y^2)^2-8xye^{x^2+y^2}$ over the unit disk. The odd
term vanishes and the first term integrates to $\pi$. This confirms the answer.}
\alternative{Split the boundary integral into
$\int z^2(-y\,\dd x+x\,\dd y)$ and
$\int(e+ze^z)\,\dd z$. The second is an exact one-variable differential,
and the first becomes $\int_0^{2\pi}\cos^2 2t\,\dd t$. This is the same
Stokes-to-boundary method in differential notation.}

\solutionproblem{7}{Operators, physical meaning, and choice of method}{14}
\textbf{(a)} In order, the answers are
\[
\boxed{\text{vector},\quad\text{scalar},\quad\text{undefined},\quad\text{scalar}}.
\]
Both $\grad f$ and $\vec F$ are vectors, so their cross product is a
vector and its curl is a vector. A dot product of two vectors is scalar.
The divergence of $\vec F$ is scalar, so taking its curl is undefined.
The curl of $\vec F$ is vector-valued, and taking its divergence is
well-defined and scalar (in fact, identically zero).

\textbf{(b)} For $\vec V=(-y,x,z)$,
\[
\dv\vec V=1>0,\qquad\curl\vec V=(0,0,2).
\]
Thus there is \textbf{local volume expansion}, and the local rotation is
\textbf{counterclockwise viewed from the positive $z$-axis}. The field
vanishes at the origin, but these derivative-based local behaviors do not.

\textbf{(c)} \textbf{All four methods apply.} The curve is smooth, closed,
and lies in $z=0$, while the field is smooth on all of $\mathbb R^3$.
Direct parametrization therefore applies. Also
\[
\vec W=\grad\left(x^2y+\frac{z^2}{2}\right),
\]
so FTLI applies. On $z=0$, the first two components form the smooth
planar field $(2xy,x^2)$, so Green's theorem applies over the unit disk.
The same upward-oriented disk is a spanning surface for Stokes' theorem.
Both the planar scalar curl and the full curl vanish. No integral value
is requested; every method would give zero.
\verification{The potential's three derivatives are $(2xy,x^2,z)$.
For the physical part, the positive $z$-component of curl follows the
right-hand rule for the planar pair $(-y,x)$.}
\alternative{For (c), checking zero curl on all of $\mathbb R^3$ also
establishes conservativity, since this domain is simply connected; the
explicit potential is the shortest justification.}

\clearpage
\setcounter{page}{1}
{\large\bfseries Mock Exam 3 --- Version C --- Solutions}
\smallskip
% Worked solutions; provenance notes are retained in comments below.
\solutionproblem{1}{Curl, divergence, and conservativity}{10}
\insight{The three rational terms have the same positive denominator.
Their mixed derivatives cancel, and the polynomial part is also a gradient.}
Put $q=x^2+y^2+z^2$, so $D=1+q$. If $\vec F=(P,Q,R)$, then
\[
P_y=Q_x=z-\frac{4xy}{D^2},\qquad
P_z=R_x=y-\frac{4xz}{D^2},\qquad
Q_z=R_y=x-\frac{4yz}{D^2}.
\]
Consequently
\[
\boxed{\curl\vec F=\vec0}.
\]
For the divergence,
\[
\begin{aligned}
\dv\vec F
&=\left(\frac2D-\frac{4x^2}{D^2}\right)
+\left(\frac2D-\frac{4y^2}{D^2}\right)
+\left(\frac2D-\frac{4z^2}{D^2}\right)\\
&=\frac6D-\frac{4q}{D^2}
=\boxed{\frac{2(q+3)}{(1+q)^2}}.
\end{aligned}
\]
The field is $C^1$ on all of $\mathbb R^3$, so its zero curl implies
conservativity here. A direct verification is
\[
\boxed{\vec F=\grad\bigl(xyz+\log(1+x^2+y^2+z^2)\bigr)}.
\]
\verification{Differentiating the displayed potential recovers all three
components. At the origin the divergence is $6$, agreeing with the
unsimplified derivative sum.}
\alternative{Recognize the potential first, then calculate the divergence
by summing its three second partial derivatives. The requested curl and
divergence still need to be supplied.}

\solutionproblem{2}{Green's theorem}{12}
\insight{The exponential terms cancel in the scalar curl. Completing
the square locates the disk; polar coordinates centered at $(1,0)$
make the remaining odd term integrate to zero.}
For $P=ye^{xy}$ and $Q=xe^{xy}+x^3/3+xy^2$,
\[
Q_x=e^{xy}+xye^{xy}+x^2+y^2,\qquad
P_y=e^{xy}+xye^{xy},
\]
so $Q_x-P_y=x^2+y^2$. Since $C$ is positively oriented and the field
is smooth on and inside $C$, Green's theorem gives
\[
\oint_C P\,\dd x+Q\,\dd y=\iint_R(x^2+y^2)\,\dd A.
\]
The region is $(x-1)^2+y^2\leq1$, $y\geq0$. Set
\[
x=1+\rho\cos\theta,\qquad y=\rho\sin\theta,\qquad
\left\{\begin{aligned}
&0\leq\theta\leq\pi,\\
&0\leq\rho\leq1.
\end{aligned}\right.
\]
The Jacobian is $\rho$, so
\[
\begin{aligned}
\iint_R(x^2+y^2)\,\dd A
&=\int_0^\pi\int_0^1(1+2\rho\cos\theta+\rho^2)\rho\,\dd\rho\,\dd\theta\\
&=\int_0^\pi\left(\frac12+\frac23\cos\theta+\frac14\right)\dd\theta
=\boxed{\frac{3\pi}{4}}.
\end{aligned}
\]
\verification{Ordinary polar coordinates instead give
$0\leq\theta\leq\pi/2$, $0\leq r\leq2\cos\theta$, and
\[
\int_0^{\pi/2}4\cos^4\theta\,\dd\theta=\frac{3\pi}{4}.
\]}
\alternative{Cartesian bounds $0\leq x\leq2$,
$0\leq y\leq\sqrt{2x-x^2}$ are correct. Any correct evaluation of
the resulting Green double integral is acceptable.}

\solutionproblem{3}{Surface parametrizations}{12}
\begin{enumerate}[label=(\alph*)]
\item The ellipsoid parametrization is
\[
\boxed{\vec r(\theta,\phi)=
(3\sin\phi\cos\theta,\ 2\sin\phi\sin\theta,\ \cos\phi)},\qquad
\left\{\begin{aligned}
&0\leq\theta\leq2\pi,\\
&0\leq\phi\leq\pi/6.
\end{aligned}\right.
\]
Indeed, $x^2/9+y^2/4=\sin^2\phi$. On the upper half,
$0\leq\phi\leq\pi/2$, the cylinder condition is
$\sin\phi\leq1/2$, hence $\phi\leq\pi/6$.

\item Use $\vec r(r,\theta)=(r\cos\theta,r\sin\theta,r)$. The first
octant initially gives $0\leq\theta\leq\pi/2$. The plane condition is
\[
r\geq r\sin\theta+1
\quad\Longleftrightarrow\quad
r\geq\frac1{1-\sin\theta}.
\]
This is compatible with $r\leq2$ exactly when
$\sin\theta\leq1/2$. Thus a complete answer is
\[
\boxed{\vec r(r,\theta)=(r\cos\theta,r\sin\theta,r)},\qquad
\left\{\begin{aligned}
&0\leq\theta\leq\pi/6,\\
&\frac1{1-\sin\theta}\leq r\leq2.
\end{aligned}\right.
\]

\item Interpolate between the two corresponding endpoints:
\[
\boxed{\vec r(s,t)=(1-t)\vec a(s)+t\vec b(s)
=(s(1-t),\ st,\ s(1-t)+2t)},
\]
with $0\leq s\leq1$, $0\leq t\leq1$.
At $t=0$ and $t=1$ the two specified lines are recovered; fixing $s$
and varying $t$ gives their joining segment.
\end{enumerate}
\verification{In (a), substitution gives the ellipsoid equation. In
(b), the lower bound is exactly the slanted plane and the upper bound
is $z=2$. In (c), recover $s=x+y$ and $t=(z-x)/2$, so the interior
is not multiply covered. Standard seam/pole coincidences on parameter
boundaries in (a) do not change the covered surface.}
\alternative{Any equivalent two-parameter representation with the
correct image and full domain is valid. For example, graph coordinates
are possible in (a), and the order of the two parameters may be reversed.}

\solutionproblem{4}{Scalar surface integral: setup only}{12}
\insight{The radial coordinate is $u>0$, so the exponential's quotient
becomes $v$. The mixed trigonometric terms in the area factor cancel
after squaring, but the parameter region is still triangular.}
Write $c=\cos v$, $s=\sin v$. Then
\[
\vec r_u=(c,s,v),\qquad \vec r_v=(-us,uc,u),
\]
and
\[
\vec r_u\times\vec r_v
=u(s-vc,-c-vs,1).
\]
Its magnitude is
\[
\begin{aligned}
\|\vec r_u\times\vec r_v\|
&=u\sqrt{(s-vc)^2+(c+vs)^2+1}\\
&=u\sqrt{2+v^2}.
\end{aligned}
\]
Since $x^2+y^2=u^2$ and $z/\sqrt{x^2+y^2}=uv/u=v$,
\[
f(\vec r(u,v))=u^2e^v.
\]
Therefore the requested setup is
\[
\boxed{\iint_S f\,\dd\sigma
=\int_1^2\int_0^{u-1}u^3e^v\sqrt{2+v^2}\,\dd v\,\dd u}.
\]
No evaluation is requested.
\verification{The cross product's third component is $u\geq1$, so the
map is regular. The radius $u$ and angle $0\leq v\leq1$ identify each
surface point uniquely. The displayed integrand is a nonnegative scalar.}
\alternative{Reversing integration order gives
$\int_0^1\int_{v+1}^2 u^3e^v\sqrt{2+v^2}\,\dd u\,\dd v$.
Reversing the cross-product order gives the same magnitude.}

\solutionproblem{5}{Direct computation of oriented flux}{20}
\insight{The cylinder is sheared in the $x$ direction as $z$ changes.
The exponential part of the field is parallel to the cylinder's ruling
direction $(1,0,1)$, so it contributes no normal flux.}
Parametrize by
\[
\vec r(\theta,z)=(z+\cos\theta,\sin\theta,z),\qquad
\left\{\begin{aligned}
&0\leq\theta\leq2\pi,\\
&0\leq z\leq2+\sin\theta.
\end{aligned}\right.
\]
The height bound is positive for every $\theta$. The tangent vectors are
\[
\vec r_\theta=(-\sin\theta,\cos\theta,0),\qquad
\vec r_z=(1,0,1),
\]
so
\[
\vec n=\vec r_\theta\times\vec r_z
=(\cos\theta,\sin\theta,-\cos\theta).
\]
For $g(x,y,z)=(x-z)^2+y^2$, this equals $\tfrac12\grad g$ on $S$,
so it points toward the exterior $g>1$. At $\theta=0$, $z=1$, for
example, the point is $(2,0,1)$ and $\vec n=(1,0,-1)$ points outward.

Put $c=\cos\theta$, $s=\sin\theta$ and $E=e^{2z+c+s}$. Substitution gives
\[
\vec F(\vec r)=(zc+E,zs,E),\qquad
\vec F(\vec r)\cdot\vec n
=zc^2+Ec+zs^2-Ec=z.
\]
Thus
\[
\begin{aligned}
\iint_S\vec F\cdot\vec N\,\dd\sigma
&=\int_0^{2\pi}\int_0^{2+\sin\theta}z\,\dd z\,\dd\theta\\
&=\frac12\int_0^{2\pi}(2+\sin\theta)^2\,\dd\theta\\
&=\frac12(8\pi+0+\pi)=\boxed{\frac{9\pi}{2}}.
\end{aligned}
\]
\verification{The exponential vector is $E(1,0,1)$, whose dot product
with $\vec n$ is zero. The remaining field has dot product $z$, agreeing
with the explicit substitution. All $z$ values are nonnegative, so the
positive final sign is also consistent.}
\alternative{A second parameter $v\in[0,1]$ may encode height through
$z=v(2+\sin\theta)$. The corresponding cross product and Jacobian must
then be included. It produces the same flux by direct parametrization.}

\solutionproblem{6}{Stokes' theorem}{20}
\insight{The diagonal edge of the base triangle lifts to a parabolic
arc, not a straight segment in the plane $z=0$. On that arc the difficult
terms form a one-variable total derivative whose endpoint values agree.}
The downward orientation induces a clockwise boundary when projected
onto the $xy$-plane. A single piecewise parametrization following
$(0,0,0)\to(0,1,0)\to(1,0,0)\to(0,0,0)$ is
\[
\vec c(t)=
\begin{cases}
(0,t,0),&0\leq t\leq1,\\
(s,1-s,s(1-s)),\quad s=t-1,&1\leq t\leq2,\\
(3-t,0,0),&2\leq t\leq3.
\end{cases}
\]
Each piece is on $S$, and the second piece is the lift of $x+y=1$.
Its derivative, piece by piece, is
\[
\vec c\,'(t)=
\begin{cases}
(0,1,0),&0<t<1,\\
(1,-1,1-2s),&1<t<2,\\
(-1,0,0),&2<t<3.
\end{cases}
\]
On the first and last pieces, respectively,
\[
\vec F(\vec c(t))=(t,0,0),\qquad
\vec F(\vec c(t))=(0,3-t,0),
\]
so both dot products are zero. On the middle piece put $w=s(1-s)$.
Then
\[
\vec F(\vec c(t))=
\bigl((1-s)(e^w+w^2),\ s(e^w+w^2)+w,\ 2w^2\bigr),
\]
and
\[
\begin{aligned}
\vec F(\vec c(t))\cdot\vec c\,'(t)
&=(1-2s)(e^w+3w^2)-w\\
&=\frac{\dd}{\dd s}(e^w+w^3)-w.
\end{aligned}
\]
Stokes' theorem now gives the required boundary line integral:
\[
\begin{aligned}
\iint_S(\curl\vec F)\cdot\vec N\,\dd\sigma
&=\oint_{\partial S}\vec F\cdot\dd\vec r\\
&=\int_0^1\left[\frac{\dd}{\dd s}(e^{s(1-s)}+s^3(1-s)^3)
-s(1-s)\right]\dd s\\
&=\bigl[e^w+w^3\bigr]_{s=0}^{s=1}
-\left[\frac{s^2}{2}-\frac{s^3}{3}\right]_0^1
=\boxed{-\frac16}.
\end{aligned}
\]
\verification{For an independent check only,
$\vec F=\grad(e^{xy}+xyz^2)+(0,z,0)$, so $\curl\vec F=(-1,0,0)$.
The downward area vector of the graph is $(y,x,-1)\,\dd x\,\dd y$.
Its dot product with this curl is $-y$, whose integral over the base
triangle is $-1/6$. }
\alternative{After identifying and parametrizing the boundary, recognize
the displayed potential decomposition. Its gradient contribution has
zero closed-curve integral, leaving $\oint z\,\dd y$, which is
$-\int_0^1s(1-s)\,\dd s$.}

\solutionproblem{7}{Conceptual questions}{14}
\begin{enumerate}[label=(\alph*)]
\item The answers are
\[
\begin{array}{c|c|c}
\text{Expression}&\text{Type}&\text{Must be zero?}\\\hline
\text{(i)}&\text{scalar}&\text{yes}\\
\text{(ii)}&\text{undefined}&\text{not applicable}\\
\text{(iii)}&\text{scalar}&\text{no}
\end{array}
\]
The field $f\vec V$ is a vector field, and the divergence of its curl
is zero. In (ii), $\dv\vec V$ is scalar, so taking its curl is not a
defined operation in the course. For (iii), $f\grad f$ is vector-valued;
its divergence need not vanish. The choice $f=x$ gives
$\dv(f\grad f)=\dv(x,0,0)=1$.

\item The correct selections are \textbf{A and B}. Direct parametrization
applies to the given smooth curve. Also
$\vec H=\grad(xyz)$, so the Fundamental Theorem of Line Integrals applies.
The curve is open: its endpoints are $(0,0,0)$ and $(1,1,1)$.
Neither Green's theorem nor Stokes' theorem applies directly to it as
stated. No integral evaluation was requested.

\item The local rotational axis is the positive $x$ direction, with
\textbf{counterclockwise} tendency viewed from the positive $x$-axis.
Zero divergence indicates \textbf{neither} local net volume expansion
nor local net volume compression. It does not force the curl to vanish.
Likewise, a zero velocity at one point does not force its spatial
derivatives to vanish.
\end{enumerate}
\verification{The local data in (c) are consistent, for example, with
$\vec v=(0,-z,y)$ at the origin: its curl is $(2,0,0)$ and
its divergence is zero.}
\alternative{Equivalent physical language such as local contraction
instead of sink is valid. Part (c) does not ask for angular speed.}

\clearpage
\setcounter{page}{1}
{\large\bfseries Version D --- Complete solutions}
\par\textit{Instructor material. Key observations and checks are included
for every problem. Sources and scope notes follow the solutions.}

\solutionproblem{1}{Curl, divergence, and global conservativity}{10}
\insight{The domain omits the whole line $x=z$, $y=0$.
Zero curl on this domain does not, by itself, settle global conservativity.}
Put $a=x-z$ and write
\[
\vec F_\lambda=\grad(xy+z^2)+\lambda\vec W,
\qquad \vec W=\left(-\frac yq,\frac aq,\frac yq\right),
\qquad q=a^2+2y^2.
\]
Here $q_x=2a$, $q_y=4y$, and $q_z=-2a$. The mixed-component
derivatives satisfy
\[
\begin{aligned}
(W_3)_y=(W_2)_z&=\frac{a^2-2y^2}{q^2},\\
(W_1)_z=(W_3)_x&=-\frac{2ay}{q^2},\\
(W_2)_x=(W_1)_y&=\frac{2y^2-a^2}{q^2}.
\end{aligned}
\]
Thus $\curl\vec W=\vec0$, and
\[
\dv\vec W=\frac{2ay-4ay+2ay}{q^2}=0.
\]
Since the base gradient has divergence $2$,
\[
\boxed{\curl\vec F_\lambda=\vec0,\qquad
\dv\vec F_\lambda=2\quad\hbox{for every }\lambda.}
\]
To decide conservativity, use the closed ellipse
\[
\vec c(t)=\left(\cos t,\frac{\sin t}{\sqrt2},0\right),
\quad0\le t\le2\pi,\qquad
\vec c\,'(t)=\left(-\sin t,\frac{\cos t}{\sqrt2},0\right).
\]
It lies entirely in $U$ because $q(\vec c(t))=1$. Direct substitution gives
\[
\vec W(\vec c(t))\cdot\vec c\,'(t)=\frac1{\sqrt2},\qquad
\oint_c\grad(xy+z^2)\cdot\dd\vec r=0.
\]
Consequently
\[
\oint_c\vec F_\lambda\cdot\dd\vec r=\lambda\sqrt2\pi.
\]
A conservative field has zero integral around every closed curve.
Hence $\lambda\ne0$ is impossible. For $\lambda=0$, the displayed
base potential works on all of $U$. The complete answer is
\[
\boxed{\vec F_\lambda\text{ is conservative on }U
\ \Longleftrightarrow\ \lambda=0,\qquad f=xy+z^2+C.}
\]
\verification{The same obstruction is visible in orthonormal coordinates
$X=(x-z)/\sqrt2$, $Y=y$, $Z=(x+z)/\sqrt2$:
\[
\vec W\cdot\dd\vec r
=\frac1{\sqrt2}\frac{-Y\,\dd X+X\,\dd Y}{X^2+Y^2}.
\]
A once-around circle in the $XY$-plane therefore has circulation
$\sqrt2\pi$, agreeing with the explicit ellipse calculation. This
coordinate observation is a check, not a needed topological theorem.}

\solutionproblem{2}{Green's theorem}{12}
\insight{Green's theorem removes the exponential. The resulting
single-variable integral is best paired with itself after a substitution.}
With $P=ye^{xy}-y/(1+x^2)$ and $Q=xe^{xy}$,
\[
Q_x=e^{xy}+xye^{xy},\qquad
P_y=e^{xy}+xye^{xy}-\frac1{1+x^2}.
\]
The field is smooth near the whole region, and the prescribed boundary
orientation is positive. Therefore
\[
\begin{aligned}
I&=\oint_C P\,\dd x+Q\,\dd y
=\int_0^1\int_0^{\ln(1+x)}\frac1{1+x^2}\,\dd y\,\dd x\\
&=\int_0^1\frac{\ln(1+x)}{1+x^2}\,\dd x.
\end{aligned}
\]
Use $x=(1-t)/(1+t)$, which exchanges $0$ and $1$. Then
\[
\frac{\dd x}{1+x^2}=-\frac{\dd t}{1+t^2},\qquad
\ln(1+x)=\ln2-\ln(1+t).
\]
Thus
\[
I=\int_0^1\frac{\ln2-\ln(1+t)}{1+t^2}\,\dd t
=\frac\pi4\ln2-I,
\]
and
\[
\boxed{I=\frac{\pi\ln2}{8}.}
\]
\verification{Independently set $x=\tan\theta$, $0\le\theta\le\pi/4$.
If $h(\theta)=\ln(1+\tan\theta)$, the tangent subtraction formula gives
$h(\theta)+h(\pi/4-\theta)=\ln2$. Integrating the pair gives the same
answer. Numerically the value is approximately $0.2721982613$.}
\alternative{The reflection in the check is an equally valid evaluation
of the double integral after Green's theorem has been used. Direct
boundary integration alone does not meet the requested method.}

\solutionproblem{3}{Three surface parametrizations}{12}
\insight{Use orthogonal projection for the tilted axis, straight rulings
for the hyperboloid, and a dependent angular bound for the sphere.}
\begin{enumerate}[label=(\alph*)]
\item Let
\[
\vec e=\frac{(1,1,0)}{\sqrt2},\qquad
\vec b=\frac{(1,-1,0)}{\sqrt2},\qquad \vec k=(0,0,1).
\]
The axis is parallel to $\vec e$, while $\vec b$ and $\vec k$ form
an orthonormal basis of its perpendicular plane. The projection of
$\vec c(t)$ onto the axis and its distance from the axis are
\[
\vec p(t)=(\vec c(t)\cdot\vec e)\vec e
=\left(\frac t2,\frac t2,0\right),\qquad
R(t)=\sqrt{\frac{t^2}{2}+t^4}.
\]
Rotating a point about the axis sweeps the entire circle with this center
and radius. Hence one answer is
\[
\boxed{\vec r(t,\theta)=\vec p(t)+R(t)
\bigl(\vec b\cos\theta+\vec k\sin\theta\bigr)},\qquad
\left\{\begin{aligned}
&1\le t\le2,\\
&0\le\theta\le2\pi.
\end{aligned}\right.
\]
Its angular starting point need not be the original generating point;
each whole circle is unchanged by a shift of angular origin. For a
parametrization with $\vec r(t,0)=\vec c(t)$ instead, one can use
\[
\vec r(t,\theta)=\vec p(t)
+\left(\frac t{\sqrt2}\cos\theta+t^2\sin\theta\right)\vec b
+\left(t^2\cos\theta-\frac t{\sqrt2}\sin\theta\right)\vec k.
\]
This rotates the perpendicular vector
$(t/\sqrt2)\vec b+t^2\vec k$ with unchanged length.

\item Start at $(\cos u,\sin u,0)$ and take a straight line with
direction $(-\sin u,\cos u,1)$. This yields
\[
\boxed{\vec r(u,v)=(\cos u-v\sin u,\ \sin u+v\cos u,\ v)},
\qquad
\left\{\begin{aligned}
&0\le u\le2\pi,\\
&0\le v\le1.
\end{aligned}\right.
\]
At fixed $u$ this is affine in $v$, as required. Moreover,
\[
x^2+y^2=(\cos u-v\sin u)^2+(\sin u+v\cos u)^2=1+v^2,
\]
so it lies on the hyperboloid. At each height $z=v$, the first two
coordinates are a rotation and dilation of the unit circle, and thus
cover the complete circle $x^2+y^2=1+v^2$. This proves full coverage.

\item Use the usual sphere parametrization
\[
\boxed{\vec r(\phi,\theta)
=(2\sin\phi\cos\theta,2\sin\phi\sin\theta,2\cos\phi)}.
\]
The extra inequality forces $z\ge0$, so $0\le\phi\le\pi/2$. It becomes
\[
\cos\phi\ge\sin\phi\bigl(|\cos\theta|+|\sin\theta|\bigr).
\]
The complete bounds are therefore
\[
\boxed{\left\{\begin{aligned}
&0\le\theta\le2\pi,\\
&0\le\phi\le\arctan\!\left(
\frac1{|\cos\theta|+|\sin\theta|}\right).
\end{aligned}\right.}
\]
The denominator is at least $1$, so this expression is defined at every
azimuth and correctly includes all four quadrants.
\end{enumerate}
\verification{For (a), $\vec r\cdot\vec e=t/\sqrt2$ and
$\|\vec r-\vec p(t)\|=R(t)$, precisely the invariants of each rotation
orbit; the axis coordinate is strictly increasing with $t$. For (b),
$\vec r_u\times\vec r_v=(\cos u-v\sin u,\sin u+v\cos u,-v)$ has squared
norm $1+2v^2>0$. For (c), substitution gives sphere radius $2$, and the
inequality was transformed equivalently, so the bounds neither omit nor
add a patch. The usual pole and angular-seam repetitions are harmless.}
\alternative{In (b), the other family of straight rulings, obtained by
reversing the tangential horizontal direction, is also valid. A generic
height-and-angle parametrization without straight parameter curves does
not meet that part's additional condition. Equivalent two-parameter
maps and equivalent complete bounds are otherwise valid.}

\solutionproblem{4}{A scalar surface integral: setup only}{12}
\insight{The map doubles the polar angle. Its full parameter annulus
covers the same conical surface twice, and $\sqrt{u^2}$ is $|u|$.}
The required derivatives are
\[
\vec r_u=(2u,2v,2u),\qquad \vec r_v=(-2v,2u,2v),
\]
\[
\vec r_u\times\vec r_v
=(4v^2-4u^2,-8uv,4u^2+4v^2),\qquad
\|\vec r_u\times\vec r_v\|=4\sqrt2(u^2+v^2).
\]
The substituted function is
\[
f(\vec r(u,v))=\sqrt{2u^2}+e^{u^2+v^2}
=\sqrt2|u|+e^{u^2+v^2}.
\]
To examine coverage, put $u=s\cos\alpha$, $v=s\sin\alpha$. Then
\[
\vec r=(s^2\cos2\alpha,s^2\sin2\alpha,s^2),\qquad1\le s\le2.
\]
The full range $0\le\alpha\le2\pi$ travels around each surface circle
twice. The half-annulus $v\ge0$, equivalently $0\le\alpha\le\pi$,
covers the geometric surface once except on its seam. Its image is
$z=\sqrt{x^2+y^2}$, $1\le z\le4$.

The polar Jacobian in the \emph{parameter plane} is $s$, in addition
to the already computed surface factor. Thus a complete unevaluated
setup is
\[
\boxed{\iint_S f\,\dd\sigma
=\int_0^\pi\int_1^2
\bigl(\sqrt2 s|\cos\alpha|+e^{s^2}\bigr)
4\sqrt2 s^3\,\dd s\,\dd\alpha.}
\]
\verification{The cross-product norm never vanishes on the annulus.
Equality of two images forces the same $s$ and angles differing by a
multiple of $\pi$, so there are exactly two preimages in the full annulus
away from seams. As a separate check, the once-around cone map
$(w\cos\beta,w\sin\beta,w)$ has area factor $\sqrt2 w$;
$w=s^2$, $\beta=2\alpha$ reproduce $4\sqrt2 s^3\,\dd s\,\dd\alpha$.}
\alternative{One-half of the full-annulus parameter integral is equivalent:
\[
\frac12\int_0^{2\pi}\int_1^2
\bigl(\sqrt2 s|\cos\alpha|+e^{s^2}\bigr)
4\sqrt2 s^3\,\dd s\,\dd\alpha.
\]
This is a correction for repeated coverage, not for orientation. A
scalar surface integral uses the cross-product magnitude.}

\solutionproblem{5}{Direct flux}{20}
\insight{The field is radial, but the upward normal changes from pointing
partly with the field to pointing partly against it. Reciprocal radii
pair these contributions exactly.}
Use the direct parametrization
\[
\vec r(r,\theta)=(r\cos\theta,r\sin\theta,1+r^2),\qquad
\left\{\begin{aligned}
&\tfrac12\le r\le2,\\
&0\le\theta\le2\pi.
\end{aligned}\right.
\]
Then
\[
\vec r_r=(\cos\theta,\sin\theta,2r),\quad
\vec r_\theta=(-r\sin\theta,r\cos\theta,0),
\]
\[
\vec n=\vec r_r\times\vec r_\theta
=(-2r^2\cos\theta,-2r^2\sin\theta,r).
\]
Its third component $r$ is positive everywhere on the parameter domain,
so it gives the required upward orientation. With $D(r)=r^4+3r^2+1$,
\[
\vec H(\vec r)=
\frac{(r\cos\theta,r\sin\theta,1+r^2)}{D(r)^{3/2}},\qquad
\vec H(\vec r)\cdot\vec n
=\frac{r(1-r^2)}{D(r)^{3/2}}.
\]
Therefore the flux is
\[
\Phi=\int_0^{2\pi}\int_{1/2}^2
\frac{r(1-r^2)}{(r^4+3r^2+1)^{3/2}}\,\dd r\,\dd\theta.
\]
One way to evaluate it is to observe that
\[
\frac{\dd}{\dd r}\frac{1+r^2}{\sqrt{r^4+3r^2+1}}
=\frac{r(r^2-1)}{(r^4+3r^2+1)^{3/2}}.
\]
The endpoint values of this primitive are both $5/\sqrt{29}$:
\[
\boxed{\Phi=-2\pi\left[
\frac{1+r^2}{\sqrt{r^4+3r^2+1}}\right]_{1/2}^2=0.}
\]
\verification{There is an independent evaluation without guessing a
primitive. Put $g(r)=r(1-r^2)/(r^4+3r^2+1)^{3/2}$. Direct algebra gives
\[
\frac{g(1/s)}{s^2}=-g(s),\qquad
\int_1^2g(r)\,\dd r=-\int_{1/2}^1g(s)\,\dd s.
\]
The two radial integrals cancel. The signed density is positive for
$r<1$ and negative for $r>1$, so zero is compatible with the field being
nonzero everywhere. Numerical integration also gives zero to numerical
precision. No divergence theorem or added surface is used.}

\solutionproblem{6}{Stokes' theorem}{20}
\insight{The sphere and cylinder intersect in a nonplanar curve. A
half-angle identity makes its height elementary and nonnegative.}
The cylindrical condition is $x^2+y^2\le2x$. Its boundary on the sphere
has
\[
x^2+y^2=2x,\qquad z^2=4-2x,\qquad z\ge0.
\]
Parametrize the projected circle by $x=1+\cos t$, $y=\sin t$.
Then
\[
z=\sqrt{2-2\cos t}=2\sin(t/2),\qquad0\le t\le2\pi,
\]
where the sine is nonnegative on this interval. Hence the complete
boundary is
\[
\vec c(t)=(1+\cos t,\sin t,2\sin(t/2)),\qquad0\le t\le2\pi.
\]
There is no additional equatorial arc: on $z=0$, the sphere has
$x^2+y^2=4$, while the cylinder inequality demands $4\le2x$;
the only possible point is $(2,0,0)$. This is already the starting and
ending point of $\vec c$.

On the rest of $S$ the outward sphere normal has positive third component.
Thus the counterclockwise projected circle gives the induced boundary
orientation. The corner of the closed boundary at $(2,0,0)$ is allowed
for a piecewise smooth boundary. Its derivative is
\[
\vec c\,'(t)=(-\sin t,\cos t,\cos(t/2)).
\]
Write $s=\sin(t/2)$, $c=\cos(t/2)$. On this curve,
\[
\vec F(\vec c(t))=(0,4s^2,-4s^2c),
\]
and
\[
\vec F(\vec c(t))\cdot\vec c\,'(t)
=4s^2(c^2-s^2)-4s^2c^2=-4s^4.
\]
Stokes' theorem therefore gives
\[
\begin{aligned}
\iint_S(\curl\vec F)\cdot\vec N\,\dd\sigma
&=\oint_{\partial S}\vec F\cdot\dd\vec r
=-4\int_0^{2\pi}\sin^4(t/2)\,\dd t\\
&=-8\int_0^\pi\sin^4u\,\dd u
=\boxed{-3\pi}.
\end{aligned}
\]
For the last integral, $\sin^4u=(3-4\cos2u+\cos4u)/8$ gives
$\int_0^\pi\sin^4u\,\dd u=3\pi/8$.
\verification{Independently, $\curl\vec F=(-3z,0,0)$. On the upper
sphere graph, the outward area vector is $(x/z,y/z,1)\,\dd x\,\dd y$
where $z>0$. The dot product is $-3x$, which extends continuously to
the lone equatorial boundary point. Over the disk centered at $(1,0)$
of radius $1$, its integral is $-3\pi$. This check agrees with Stokes;
the required solution above actually uses the boundary line integral.}

\solutionproblem{7}{Operators, physical meaning, and a global obstruction}{14}
\insight{Multiplying a curl by a variable scalar can create divergence.
Also, a global vector potential requires more than local zero divergence
on a domain with an excluded point.}
\begin{enumerate}[label=(\alph*)]
\item The expression is a defined \textbf{scalar field}. The divergence
product rule and the identity $\dv(\curl\vec G)=0$ give
\[
\boxed{\dv(f\,\curl\vec G)=\grad f\cdot(\curl\vec G).}
\]
It need not vanish. For example,
\[
f=z,\qquad \vec G=(-y/2,x/2,0),\qquad
\curl\vec G=(0,0,1),
\]
so $\dv(f\,\curl\vec G)=\dv(0,0,z)=1$.

\item Put $\rho=\sqrt{x^2+y^2+z^2}>0$. Direct differentiation gives
\[
\dv\vec H
=\left(\rho^{-3}-3x^2\rho^{-5}\right)
+\left(\rho^{-3}-3y^2\rho^{-5}\right)
+\left(\rho^{-3}-3z^2\rho^{-5}\right)=0.
\]
Also $\vec H=\grad(-1/\rho)$, so $\curl\vec H=\vec0$ on $U_0$.
Thus at every point of $U_0$ these operators predict \textbf{neither
local volume expansion nor compression, and no local rotational tendency}.
They do not say that the velocity is zero or that the motion cannot stretch
in some directions while compressing in others.

For a sphere of radius $R$ about the origin, however,
\[
\vec H\cdot\vec N=\frac1{R^2},\qquad
\iint_{\rho=R}\vec H\cdot\vec N\,\dd\sigma=4\pi.
\]
This outward flux encloses the excluded origin. No conclusion of zero
flux through such a sphere follows just from divergence zero at every
point where the field is defined. There is no contradiction with the
local source-free interpretation away from the origin.

\item \textbf{No such global $\vec A$ exists.} Suppose it did. Split the
unit sphere into its upper and lower hemispheres $S_+$ and $S_-$, both
oriented outward. Their common equator is counterclockwise when induced
by $S_+$ and clockwise when induced by $S_-$. Since $\vec A$ would be
smooth near both hemispheres, Stokes' theorem would imply
\[
\begin{aligned}
\iint_{S^2}\vec H\cdot\vec N\,\dd\sigma
&=\iint_{S_+}\curl\vec A\cdot\vec N\,\dd\sigma
+\iint_{S_-}\curl\vec A\cdot\vec N\,\dd\sigma\\
&=\oint_{C}\vec A\cdot\dd\vec r
+\oint_{-C}\vec A\cdot\dd\vec r=0.
\end{aligned}
\]
But the actual flux, computed in (b), is $4\pi$. This contradiction
rules out a $C^2$ vector potential defined on all of $U_0$. The argument
uses only Stokes on two ordinary smooth hemispheres and cancellation of
their common oriented boundary.
\end{enumerate}
\verification{For the unit sphere the flux calculation can be checked
directly with
$\vec r(\phi,\theta)=(\sin\phi\cos\theta,\sin\phi\sin\theta,\cos\phi)$:
the outward cross product $\vec r_\phi\times\vec r_\theta
=\sin\phi\,\vec r$ gives
$\int_0^{2\pi}\int_0^\pi\sin\phi\,\dd\phi\,\dd\theta=4\pi$.
This proves the necessary obstruction without a theorem about domains
or a divergence-theorem argument.}

% ================= INSTRUCTOR-ONLY PROVENANCE / SCOPE =================
\clearpage
{\large\bfseries Version D --- Instructor notes}
\par\textit{Not part of the student-facing questions.}

\textbf{Purpose and alignment.}
This is an untimed challenge set, deliberately above the expected actual
Exam 3 level. It preserves the outline's seven sections and nominal
weights $10+12+12+12+20+20+14=100$, while following the revised compact
A--C practice format without displayed point counts or answer boxes.
Problem 3 has three two-parameter maps; Problem 4 supplies its map and
asks only for a setup; Problem 5 requires a direct oriented-flux
calculation; Problem 6 requires Stokes and an actual one-parameter
boundary integral. Green's theorem is required in Problem 2. For Problem
7, a possible division is $4+4+6$. No timing claim is made.

\textbf{What exceeds normal expectations.}
\begin{enumerate}[label=\textbf{\arabic*.},itemsep=5pt]
\item The zero-curl test must be distinguished from global conservativity
on a domain missing a tilted line. Students must construct an obstructing
closed curve. No fundamental-group or topology theorem is assumed.
\item The final integral requires discovery of a fractional-linear
substitution, or its trigonometric-reflection equivalent. This is a
contest-level integration observation after an ordinary Green setup.
\item Rotation about an oblique axis requires constructing circle centers
and radii. The hyperboloid must be recognized as ruled, not merely
parametrized in the most familiar way. The sphere has a genuinely
azimuth-dependent angular bound. Each uses elementary vectors and algebra.
\item The supplied regular map is not one-to-one on its full domain.
Correcting multiplicity and preserving an absolute value exceed the
routine surface-integral setups in the outline. The prompt explicitly
requires counting the geometric image once.
\item The signed radial flux cancels under reciprocal radii. The insight
is in the evaluation, not a large field or lengthy cross product.
No divergence theorem or solid-angle theorem is assumed.
\item The boundary is a sphere--cylinder intersection with a corner, not
a standard horizontal circle. Discovering the half-angle parametrization
and ruling out an extra equatorial component are the extensions.
\item The product identity itself is familiar. The contrast between
local zero divergence and a global obstruction to being a curl is beyond
ordinary Exam 3 practice. Two applications of the already taught Stokes
theorem supply the complete argument; no cohomology result is needed.
\end{enumerate}

\textbf{Origins and adaptations.}
The following notes distinguish specific mathematical seeds from original
constructions. None of the student prompts is copied from a contest.

\begin{itemize}[itemsep=5pt]
\item \textbf{Problem 1: original construction.} The tilted-line field
was constructed so its curl vanishes and its divergence is constant for
every $\lambda$, while a closed ellipse detects its nonzero circulation.
It uses the standard circulation obstruction also underlying the source
listed for Problem 7.
\item \textbf{Problem 2: Putnam 2005, A5.}
\url{https://kskedlaya.org/putnam-archive/2005.pdf}
The final one-variable logarithmic integral and the involutive substitution
are the specific borrowed seed. The region, vector field, Green's-theorem
formulation, and required intermediate setup are an original adaptation.
\item \textbf{Problems 3--4: original constructions.} The chosen oblique
revolution, ruled-hyperboloid requirement, sphere patch, and double-covered
conical annulus were constructed for this version. The complex-square
coordinate identity in Problem 4 is used only as elementary real algebra;
students need no complex analysis.
\item \textbf{Problem 5: strategy inspired by Putnam 1985, B5.}
\url{https://kskedlaya.org/putnam-archive/1985.pdf}
The borrowed strategy is reciprocal-variable pairing. The paraboloid,
annular bounds, inverse-square field, and direct flux calculation are
new. None of that source's Gaussian-integral machinery is used.
\item \textbf{Problem 6: Viviani-curve geometry.} Oliver Knill,
Harvard Math 136, Unit 1 homework, Problem 1.1:
\url{https://people.math.harvard.edu/~knill/teaching/math136/handouts/homework01.pdf}
The sphere--cylinder intersection is the geometric seed, rescaled and
restricted to its upper loop. The field, flux question, boundary
orientation task, and Stokes calculation are new. Curvature and torsion
from that source are irrelevant here and are not assumed.
\item \textbf{Problem 7(c): inspired by Putnam 1987, A5.}
\url{https://kskedlaya.org/putnam-archive/1987.pdf}
That problem uses a nonzero-circulation obstruction to a curl-free
extension. Here the analogous obstruction is nonzero sphere flux,
excluding a global vector potential. The proof uses two hemispheres
and ordinary Stokes. Parts (a)--(b) and their linkage with Problem 5
are original.
\end{itemize}

\textbf{Other research consulted, not directly adapted.}
The official IMC 2026 Problem 7 materials
(\url{https://www.imc-math.org.uk/?item=prob7s&section=problems&year=2026})
were considered for overlap and ruled-region ideas; its inequality is not
used. The official ETS \emph{GRE Mathematics Practice Book}, form GR3768,
including the symmetry integral in Question 63, was checked as a comparison
for observation-based calculus:
\url{https://www.ets.org/content/dam/ets-org/pdfs/gre/practice-book-math.pdf}.
The official IMO 1961 set
(\url{https://www.imo-official.org/assets/documents/problems/1961/1961_eng.pdf})
and Kyoto University's 2018 master's entrance basic-mathematics paper,
especially its first improper-integral problem
(\url{https://www.math.kyoto-u.ac.jp/files/master_exams/2017math_kiso.pdf}),
were also considered. Their original questions were not inserted: the
selected tasks keep the Exam 3 vector-calculus topics central.

\textbf{Verification record.}
Each solution was recomputed independently by a second reader. Checks
include the explicit circulation obstruction in Problem 1, two evaluations
of Problem 2, algebra and coverage tests for all three maps in Problem 3,
an independent cone parametrization for Problem 4, both a primitive and
reciprocal pairing for Problem 5, a direct surface check of Problem 6,
and direct sphere flux plus two-hemisphere Stokes for Problem 7.
Numerical quadrature independently agrees with
$\pi\ln2/8$, $0$, and $-3\pi$ in Problems 2, 5, and 6.
The mathematical content of Versions A--C and the original outline is unchanged.
% Instructor-only provenance retained from the original A--C solution sources.
% Version A
% \source{Original field design using the curl/divergence and
% conservativity tasks in \nolinkurl{261II2026/HW/789/HW7_withSpace.tex},
% Problems 1--8, and the product/cancellation ideas in
% \nolinkurl{261II2026/Q/7/q7-261-II2026.tex}.}
% \source{Original off-center-hole and translated-moment construction.
% The cancellation of mixed exponential derivatives is inspired by
% \nolinkurl{261II2026/Q/7/q7-261-II2026.tex}, Problem 4. Multiple positively
% oriented boundary components are covered in
% \nolinkurl{261II2026/Q/9/BaohuaQ9Review.tex}, Problem 6, and the corresponding
% boundary-orientation questions in HW9.}
% \source{Original combinations of the five parametrization families
% taught in class. Polar descriptions of shifted circular regions build
% on \nolinkurl{261II2026/Q/5/q5-261-II2026.tex}; scaled spherical
% parametrizations and generating-curve revolutions build on
% \nolinkurl{261II2026/Q/8/q8-261-II2026.tex} and HW8, Problems 1 and 5--9.
% Part (c) needs only the supplied curve and the taught rotation formula;
% no facts about a torus are assumed.}
% \source{Original implicit-relation cancellation on a two-parameter
% patch. The task and required steps follow
% \nolinkurl{261II2026/HW/789/HW8_withSpace.tex}, Problems 3--4, and
% \nolinkurl{261II2026/Merger/e33.pdf}, Problems 14--15.}
% \source{Original tangent-field cancellation construction. The direct
% flux procedure is the one in HW9, Problems 2 and 6; the geometry and
% outward/downward orientation align with the revised
% \nolinkurl{261II2026/Q/9/q9-261-II2026.tex}, Problem 2.}
% \source{Original field and numerical combination based on the two-rim
% Stokes geometry in \nolinkurl{261II2026/HW/789/HW9_withSpace.tex},
% Problems 16--17, and the course's right-hand-rule treatment of surface
% boundaries. No additional theorem is required.}
% \source{Original conceptual examples aligned with
% \nolinkurl{261II2026/HW/789/HW9_withSpace.tex}, Problem 19, and
% \nolinkurl{261II2026/Merger/e33.pdf}, Problems 37--41. The distinction
% between curl on a curve and a conservative field on an open region uses
% only the usual FTLI hypotheses; no topology counterexample is needed.}
% Version B
% \source{Original construction based on mixed-partial cancellation
% and the curl and divergence computations in HW7, Problems 1--8. No outside
% contest problem is used.}
% \source{Original field and curved region. The missing-segment strategy
% is adapted from Quiz 7 (2026), Problem 4; this version adds a nonconstant
% Green integrand, a clockwise closure, and a diagonal correction.}
% \source{Original intersection-bound variations of HW8, Problem 1, and
% the common-surface parametrizations in Merger/e33.pdf, Problem 36.}
% \source{Original patch and scalar function, using the setup format with
% a supplied parametrization in HW8, Problems 3--4. The inverse
% coordinate identities are the intended observation.}
% \source{Original synthesis of the graph-flux procedure in
% Merger/e33.pdf, Problems 27--28, with a symmetric parabolic lens.
% The antisymmetric cancellation and downward orientation are new.}
% \source{Original nonplanar graph and field. The strategy of simplifying
% a complicated field on the boundary is inspired by Quiz 9 (2025),
% Problem 4. The boundary is now nonplanar and the remaining $\dd z$ term
% requires a one-variable exact-derivative observation.}
% \source{Original scenarios aligned with HW9, Problem 19, and
% Merger/e33.pdf, Problems 38--41. The zero velocity at the origin is
% deliberately contrasted with nonzero local rotation and expansion.}
% Version C
% \source{Original construction combining the potential-function tasks
% in 2026 HW7, Problems 9--13, with the curl/divergence work in Problems
% 1--8. The rational radial gradient is new; no domain or topology
% exception is needed because $D>0$ everywhere.}
% \source{Original combination of the off-center circular-region ideas
% in the 2026 polar-coordinate materials and the Green calculations in
% HW7, Problems 14--25. The translation and the cancellation in the field
% are newly constructed rather than a numerical variant of an old question.}
% \source{Original cropped surfaces using the ellipsoid, cone, and ruled
% surface methods taught for Quiz 8. HW8, Problem 1(e), supplies the basic
% cone representation; 2025 Quiz 9, Problem 2, supplies the interpolation
% idea. The elliptical cylinder, slanted cone cut, and skew-line pairing
% are new geometric constraints.}
% \source{Original given-parametrization setup in the style of HW8,
% Problems 3--4, and the Merger practice set e33, Problems 14--15.
% The radius-dependent height and nonrectangular parameter domain are new.}
% \source{Original extension of the cylinder-under-a-sloping-plane
% geometry in HW8, Problem 8. The cylinder is now sheared, and the added
% tangent component is constructed to vanish from flux. This uses only
% ordinary surface parametrization and a dot product.}
% \source{Original combination of graph surfaces from HW8, Problem 1(d),
% the boundary conversion in HW9, Problems 13--18, and the conservative
% field ideas in HW7. The total-derivative identity follows from the
% ordinary single-variable chain rule, with no additional theorem.}
% \source{Original conceptual combinations based on HW9, Problem 19,
% the product-rule discussion accompanying 2026 Quiz 7, and Merger e33,
% Problems 37--41. The local-data example separates velocity from its
% derivatives, using only the taught curl/divergence interpretations.}
\end{document}
